% ==================================================================
% id:       lim-007
% topic:    גבולות ורציפות של פונקציות
% subtopic: גבולות חד־צדדיים
% solution: full
% ==================================================================
\textbf{
הוכיחו או הפריכו
כל אחד מהסעיפים הבאים:}
\begin{enumerate}
    \item
אם
$\displaystyle\lim_{x\to 1^+}f(x)=3$,
    אזי קיימת סביבה נקובה של
    $x=1$
    אשר בה
    $f(x)<4$.
\ifsol
\hint{שימו לב שסביבה נקובה של $1$ כוללת את שני הצדדים שלה, ובדקו מה הנתון אומר על כל צד. אם נראה לכם שהטענה לא נכונה, חפשו דוגמה נגדית פשוטה.}
\begin{solbox}
הטענה אינה נכונה.
ניתן כמה דוגמאות נגדיות.

שימו לב כי ידוע הגבול של
$f$
ב-
$1$
מימין, לכן היא מוגדרת בסביבה ימנית נקובה של
$1$,
ולא בהכרח מוגדרת בסביבה שמאלית. לדוגמא
$$.f(x)=\sqrt{x-1}+3$$

עוד דוגמא:
$$,f(x) = \case{cc}{
       5, &  x<1\\ 3, & x\geq 1}$$
השלימו את הפרטים.
\end{solbox}
\fi

    \item
אם
$\displaystyle\lim_{x\to -5^-}f(x)=-2$
    ובנוסף
    $\displaystyle\lim_{x\to -5^+}f(x)$
קיים,
    אזי קיימת סביבה נקובה של
    $x=-5$
    אשר בה
    $f(x)>-3$.
\ifsol
\hint{בדקו מה כל אחד מהגבולות החד-צדדיים אומר על ערכי $f$ בצד שלו. אם נראה לכם שהטענה לא נכונה, חפשו דוגמה נגדית פשוטה.}
\begin{solbox}
הטענה אינה נכונה.
ניקח
$$,f(x) = \case{cc}{
       -2 &:  x<-5\\
       -4 &: x\geq -5}$$
       השלמת הפרטים מוטלת על הקורא.
שימו לב כי קיימת סביבה נקובה שמאלית של
$-5$
אשר בה
$f(x)>-3$.
\end{solbox}
\fi

    \item נתון כי
    $\displaystyle\lim_{x\to2^-}f(x)$
    ו-
    $\displaystyle\lim_{x\to2^+}f(x)$
    קיימים,
    אזי קיימת סביבה נקובה של
    $x=2$
    אשר
    $f$
    חסומה
    בה.
\ifsol
\hint{בדקו מה הגדרת הגבול החד-צדדי (עם $\varepsilon$ קבוע) אומרת על ערכי $f$ בכל צד. אם נראה לכם שהטענה לא נכונה, חפשו דוגמה נגדית.}
\begin{solbox}
הטענה נכונה. נסמן
$L_1=\displaystyle\lim_{x\to2^-}f(x)$,
$L_2=\displaystyle\lim_{x\to2^+}f(x)$.
מקיום הגבול
הראשון
קיימת
$\delta_1>0$
כך ש-
$$.|f(x)-L_1|<1,\ \forall x\in(2-\delta_1,2)$$
לכן
$|f(x)|\leq 1+L_1$
לכל
$x\in(2-\delta_1,2)$.
בדומה,
מקיום הגבול
השני
קיימת
$\delta_2>0$
כך ש-
$$.|f(x)-L_2|<1,\ \forall x\in(2,2+\delta_2)$$
לכן
$|f(x)|\leq 1+L_2$
לכל
$x\in(2,2+\delta_2)$.
תהא
$\delta=\min\{\delta_1,\delta_2\}$.
לכן, לכל
$x$
המקיים
$x\in(2-\delta,2+\delta)$,
כך ש-
$x\neq 2$,
מתקיים
$$.|f(x)|\leq\max\{1+|L_1|,1+|L_2|\}$$
\end{solbox}
\fi

    \item
אם    קיים הגבול
    $\displaystyle\lim_{x\to 3}f(x)$,
    אזי קיים הגבול
    $\displaystyle\lim_{x\to 3}\lfloor f(x)\rfloor$.
\ifsol
\hint{בדקו כיצד מתנהגת $\lfloor t\rfloor$ כאשר $t$ קרוב לערך הגבול, ונסו את הטענה על פונקציות פשוטות.}
\begin{solbox}
הטענה אינה נכונה.
ניקח
$f(x)=x$.
בבירור
$\displaystyle\lim_{x\to3}f(x)=3$,
אך
$\displaystyle\lim_{x\to 3}\lfloor f(x)\rfloor$
אינו קיים.
לכל
$3<x<4$
מתקיים
$\lfloor f(x)\rfloor=3$,
לכן
$$.\ds{\lim_{x\to3^+}\lfloor f(x)\rfloor}=3$$
לכל
$2<x<3$
מתקיים
$\lfloor f(x)\rfloor=2$,
לכן
$$.\ds{\lim_{x\to3^-}\lfloor f(x)\rfloor}=2$$
לכן
$\ds{\lim_{x\to3}\lfloor f(x)\rfloor}$
לא קיים.
\end{solbox}
\fi

\item
אם
$\displaystyle\lim_{x\to -2}f(x)=\frac{1}{3}$,
    אזי קיים הגבול
    $\displaystyle\lim_{x\to -2}\lfloor f(x)\rfloor$.
\ifsol
\hint{בדקו מה הגבול אומר על ערכי $f$ בסביבה נקובה של $-2$, ומה נובע מכך על $\lfloor f(x)\rfloor$.}
\begin{solbox}
הטענה נכונה.
מכיוון ש-
$\displaystyle\lim_{x\to-2}f(x)=\tfrac{1}{3}$,
קיימת סביבה נקובה של
$(-2)$
שבה
$0<f(x)<1$
(מדוע? השלימו),
ולכן לכל
$x$
בסביבה זו
$\lfloor f(x)\rfloor=0$,
מכאן ש-
$\displaystyle\lim_{x\to-2}\lfloor f(x)\rfloor=0$.
\end{solbox}
\fi

    \item
    אם
     קיים הגבול
    $\displaystyle\lim_{x\to x_0}\left(f(x)+g(x)\right)$,
    אזי קיימים הגבולות
    $\displaystyle\lim_{x\to x_0} f(x)$
    ו-
    $\displaystyle\lim_{x\to x_0} g(x)$.
\ifsol
\hint{אם נראה לכם שהטענה נכונה, נסו להביע את $f$ בעזרת $f+g$ ו-$g$. אחרת, חפשו דוגמה נגדית פשוטה.}
\begin{solbox}
הטענה אינה נכונה.
תהיינה
$$f(x)=\frac{1}{x},\quad
g(x)=-f(x)$$
ו-
$x_0=0$.
מתקיים
$\displaystyle\lim_{x\to0}(f(x)+g(x))=0$
כי
$f(x)+g(x)=0$
לכל
$x\neq0$, ומצד שני
בבירור לא קיימים הגבולות
$\displaystyle\lim_{x\to0}f(x)$
ו-
$\displaystyle\lim_{x\to0}(g(x))$.

שימו לב כי כל עוד
$\ds{\lim_{x\to x_0}}f(x)$
לא קיים, הדוגמא
$g=-f$
"תעבוד".
\end{solbox}
\fi
\end{enumerate}
